\section{超越基}\label{sec:transcendence-bases}

如第~\ref{sec:motivation-for-dimension}~节所述, 仿射簇 $X$ 的维数应当等于 $K[X]$ 中代数独立元素的最大个数. 事实上, 在域中研究代数相关与代数独立更为直接. 这就是第~\ref{sec:function-fields}~节为不可约仿射簇引入相应的函数域 $K(X)$ 的原因. 对于不可约仿射簇, 有一组自然的包含关系
\[
  K\subseteq K[X]\subseteq K(X),
\]
且 $K[X]$ 中代数独立的元素在 $K(X)$ 中仍然代数独立. 以此为动机, 我们转向研究形如 $K\subseteq K(X)$ 的域扩张中的代数相关与代数独立.

为便于讨论域扩张, 先引入域的生成元这一概念.

\begin{Definition}{域的生成元}\leavevmode\label{def:field-generators}
设 $K\subseteq L$ 为域, 且 $S\subseteq L$ 为子集. 由 $S$ 生成的 $K$ 的\emph{域扩张}是
\[
  K(S)=\{ab^{-1}:a,b\in K[S],\ b\neq0\}\subseteq L.
\]
若对某个有限集 $S\subseteq L$ 有 $L=K(S)$, 则称 $L$ 是 $K$ 的\emph{有限生成域扩张}.
\end{Definition}

注意 $K(S)$ 是 $L$ 的同时包含 $K$ 和 $S$ 的最小子域 (Exercise~\ref{ex:generated-field-minimality}), 且它与 $\operatorname{Frac}(K[S])$ 典范同构, 其中 $K[S]$ 是由 $S$ 生成的 $K$-代数. 此外, 有 (Exercise~\ref{ex:generated-field-in-stages})
\[
  K(S_{1})(S_{2})=K(S_{1}\cup S_{2}).
\]
当 $S=\{a_{1},\ldots,a_{n}\}$ 时, 省去集合的花括号, 写成 $K(S)=K(a_{1},\ldots,a_{n})$.

\begin{Example}{有理函数域}\leavevmode\label{ex:rational-function-field-generators}
就这里的目的而言, $K$ 的域扩张的典型例子是第~\ref{sec:unique-factorization-polynomial-rings}~节引入的有理函数域:
\[
  K\subseteq K(x_{1},\ldots,x_{n}).
\]
如记号所示, $K(x_{1},\ldots,x_{n})$ 在 $K$ 上由 $x_{1},\ldots,x_{n}$ 有限生成.
\end{Example}

\begin{Example}{函数域}\leavevmode\label{ex:function-field-generated-by-coordinate-functions}
设 $X\subseteq\mathbb{A}^{n}$ 为不可约仿射簇, 坐标环为 $K[X]$, 函数域为 $K(X)$. 坐标函数 $[x_{1}],\ldots,[x_{n}]\in K[X]$ 可以自然地视为 $K(X)$ 的元素, 而 $K(X)$ 的任意元素都能写成
\[
  \frac{[f(x_{1},\ldots,x_{n})]}{[g(x_{1},\ldots,x_{n})]}
  =\frac{f([x_{1}],\ldots,[x_{n}])}{g([x_{1}],\ldots,[x_{n}])},
\]
其中 $f,g\in K[x_{1},\ldots,x_{n}]$, 且 $[g]\neq0$. 按定义, 最后一个表达式属于 $K([x_{1}],\ldots,[x_{n}])$, 这意味着 $K(X)$ 由 $[x_{1}],\ldots,[x_{n}]$ 有限生成.
\end{Example}

现在来定义超越基这一关键概念. 它是向量空间的基在域扩张中的对应概念. 回忆一下, 向量空间的基须满足两个条件: (i) 线性独立, (ii) 张成向量空间. 第一个条件给基的大小设定了上界: 一个向量空间子集若有太多元素, 它们就会线性相关. 第二个条件给基的大小设定了下界: 元素若太少, 就张不成整个向量空间. 同样, 超越基也必须满足两个性质. 如下一节所示, 第一个性质给出超越基大小的上界, 第二个则给出下界.

\begin{Definition}{超越基}\leavevmode\label{def:transcendence-basis}
设 $K\subseteq L$ 为域. 若子集 $S\subseteq L$ 满足下列两个条件, 则称它是 $L$ 在 $K$ 上的\emph{超越基}:
\begin{enumerate}
  \item[(i)] $S$ 在 $K$ 上代数独立;
  \item[(ii)] $L$ 在 $K(S)$ 上代数.
\end{enumerate}
\end{Definition}

\begin{BookFlowNote}[l]{%
在讨论超越基的具体例子之前, 先熟悉代数域扩张的一些性质会有所帮助. Definition~\ref{def:transcendence-basis} 的条件 (ii) 中就出现了这样的扩张.
}
\noindent 消去分母后可知, (ii) 与一个看似更强的条件等价, 即 $L$ 在 $K$-代数 $K[S]$ 上代数 (Exercise~\ref{ex:algebraic-over-generated-field-iff-algebraic-over-algebra}).

下面的结论是 Theorem~\ref{thm:finite-generation-integrality} 在域中的对应结论, 也是我们所需的关键工具.
\end{BookFlowNote}

\begin{Proposition}{用代数元扩张域}\leavevmode\label{prop:extending-fields-by-algebraic-elements}
设 $K\subseteq L$ 为域, 且 $a_{1},\ldots,a_{n}\in L$. 下列条件等价:
\begin{enumerate}
  \item[(i)] 对每个 $i=1,\ldots,n$, $a_{i}$ 在 $K$ 上代数;
  \item[(ii)] $K[a_{1},\ldots,a_{n}]$ 是 $K$ 上的有限维向量空间;
  \item[(iii)] $K[a_{1},\ldots,a_{n}]$ 在 $K$ 上代数;
  \item[(iv)] $K[a_{1},\ldots,a_{n}]=K(a_{1},\ldots,a_{n})$.
\end{enumerate}
\end{Proposition}
\begin{Proof}
条件 (i)--(iii) 的等价性是 Theorem~\ref{thm:finite-generation-integrality} 中 $R=K$ 为域的特殊情形. 此时代数性等价于整性, 有限生成模就是有限维向量空间.

为证明 (iv) 与其余三个条件等价, 只需证明 (iv) 与 (iii) 等价. (iv) 蕴含 (iii) 实质上就是 Zariski 引理 (Lemma~\ref{lem:zariski}) 的陈述. 具体来说, 若 $K[a_{1},\ldots,a_{n}]=K(a_{1},\ldots,a_{n})$, 则 $K[a_{1},\ldots,a_{n}]$ 是一个包含 $K$ 且作为 $K$-代数有限生成的域. 由 Zariski 引理, $K[a_{1},\ldots,a_{n}]$ 在 $K$ 上代数.

还须证明 (iii) 蕴含 (iv). 假设 $K[a_{1},\ldots,a_{n}]$ 在 $K$ 上代数. 我们要证明 $K[a_{1},\ldots,a_{n}]=K(a_{1},\ldots,a_{n})$. 包含关系 $K[a_{1},\ldots,a_{n}]\subseteq K(a_{1},\ldots,a_{n})$ 直接来自定义, 所以转而证明反向包含.

由于 $K(a_{1},\ldots,a_{n})$ 是 $L$ 中包含 $K$ 和 $a_{1},\ldots,a_{n}$ 的最小子域, 只要说明环 $K[a_{1},\ldots,a_{n}]$ 实际上是域, 就能证明 $K(a_{1},\ldots,a_{n})\subseteq K[a_{1},\ldots,a_{n}]$. 为此, 设 $a\in K[a_{1},\ldots,a_{n}]$ 非零, 必须证明 $a$ 有乘法逆元. 由 $K[a_{1},\ldots,a_{n}]$ 在 $K$ 上代数的假设, 知道存在下列形式的非平凡代数关系:
\[
  c_{d}a^{d}+\cdots+c_{1}a+c_{0}=0
    \quad\Longrightarrow\quad c_{d}a^{d}+\cdots+c_{i+1}a^{i+1}+c_{i}a^{i}=0,
\]
其中 $c_{0},\ldots,c_{d}\in K$, $c_{i}$ 是第一个非零系数. 解出 $a^{i}$, 得
\[
  a^{i}=-c_{i}^{-1}(c_{d}a^{d}+\cdots+c_{i+1}a^{i+1}).
\]
两边约去因子 $a^{i}$ (运算在域 $L$ 中进行, 所以可以这样做), 再从右边提出 $a$, 便得到 $a$ 有逆元:
\[
  1=-\bigl[c_{i}^{-1}(c_{d}a^{d-i-1}+\cdots+c_{i+1})\bigr]a.
\]
所以 $K[a_{1},\ldots,a_{n}]$ 是一个域, 如所需.
\end{Proof}

Proposition~\ref{prop:extending-fields-by-algebraic-elements} 中唯一不是 Theorem~\ref{thm:finite-generation-integrality} 特例的条件是 (iv). 用下面的例子来说明它.

\begin{Example}{{\textmd{$\mathbb{Q}(\sqrt{2})=\mathbb{Q}[\sqrt{2}]$}}}\leavevmode\label{ex:sqrt-two-field-equals-algebra}
考虑 $\sqrt{2}\in\mathbb{R}$. 由于 $\sqrt{2}$ 是多项式 $x^{2}-2\in\mathbb{Q}[x]$ 的根, 所以 $\sqrt{2}$ 在 $\mathbb{Q}$ 上代数. 因而 Proposition~\ref{prop:extending-fields-by-algebraic-elements} 断言 $\mathbb{Q}(\sqrt{2})=\mathbb{Q}[\sqrt{2}]$. 下面验证这一点.

用 Example~\ref{ex:square-root-two-finite-module} 的论证算出 $\mathbb{Q}[\sqrt{2}]=\mathbb{Q}\{1,\sqrt{2}\}$. 由域的生成元的定义, $\mathbb{Q}(\sqrt{2})$ 的每个元素都形如
\[
  \frac{a+b\sqrt{2}}{c+d\sqrt{2}},
\]
其中 $a,b,c,d\in\mathbb{Q}$, 且 $c$ 和 $d$ 不全为零. 可以将该商的分母有理化:
\[
  \frac{a+b\sqrt{2}}{c+d\sqrt{2}}
  =\frac{(a+b\sqrt{2})(c-d\sqrt{2})}{(c+d\sqrt{2})(c-d\sqrt{2})}
  =\frac{(ac-2bd)+(bc-ad)\sqrt{2}}{c^{2}-2d^{2}}=r+q\sqrt{2},
\]
其中
\[
  r=\frac{ac-2bd}{c^{2}-2d^{2}}\in\mathbb{Q}
  \qquad\text{且}\qquad
  q=\frac{bc-ad}{c^{2}-2d^{2}}\in\mathbb{Q}.
\]
这表明 $\mathbb{Q}(\sqrt{2})$ 的每个元素实际上都属于 $\mathbb{Q}[\sqrt{2}]$. 反向包含由定义得到, 故 $\mathbb{Q}(\sqrt{2})=\mathbb{Q}[\sqrt{2}]$.
\end{Example}

Proposition~\ref{prop:extending-fields-by-algebraic-elements} 的一个最重要的推论如下. 它在环的情形下不成立 (见下文 Example~\ref{ex:algebraic-transitivity-fails-for-rings}). 这正是代数相关与代数独立在域中的性质比在环中好得多的根本原因.

\begin{Corollary}{代数扩张的传递性}\leavevmode\label{cor:algebraic-field-extensions-transitive}
设 $K\subseteq L\subseteq M$ 为有限生成域扩张. 若 $L$ 在 $K$ 上代数, 且 $M$ 在 $L$ 上代数, 则 $M$ 在 $K$ 上代数.
\end{Corollary}
\begin{Proof}
由 Proposition~\ref{prop:extending-fields-by-algebraic-elements}, 该断言等价于证明: 若 $L$ 是 $K$ 上的有限维向量空间, 且 $M$ 是 $L$ 上的有限维向量空间, 则 $M$ 是 $K$ 上的有限维向量空间. 只需注意 (见 Exercise~\ref{ex:finite-module-in-tower}): 若 $L=K\{a_{1},\ldots,a_{m}\}$ 且 $M=L\{b_{1},\ldots,b_{n}\}$, 则
\[
  M=K\{a_{i}b_{j}\mid 1\leq i\leq m,\ 1\leq j\leq n\}.
\]
\end{Proof}

下一个例子说明 Corollary~\ref{cor:algebraic-field-extensions-transitive} 在环的情形下不成立.

\begin{Example}{环扩张不满足 Corollary~\ref{cor:algebraic-field-extensions-transitive}}\leavevmode\label{ex:algebraic-transitivity-fails-for-rings}
考虑环扩张
\[
  \mathbb{C}\subseteq\frac{\mathbb{C}[x]}{\langle x^{2}-1\rangle}
    \subseteq\frac{\mathbb{C}[x,y]}{\langle x^{2}-1,xy-y\rangle}.
\]
我们说明第二个环在第一个环上代数, 第三个环在第二个环上代数, 但第三个环在第一个环上不代数.

先看第一个扩张. 由关系 $[x^{2}]=1$, $\mathbb{C}[x]/\langle x^{2}-1\rangle$ 的任意元素都可以写为 $[ax+b]$, 其中 $a,b\in\mathbb{C}$. 直接计算可知, 这样的元素 $[ax+b]$ 是多项式
\[
  f(z)=z^{2}-2bz+b^{2}-a^{2}\in\mathbb{C}[z]
\]
的根, 所以第一个扩张是代数扩张.

为说明第二个扩张是代数扩张, 注意由关系 $[x^{2}]=1$ 和 $[xy]=[y]$, $\mathbb{C}[x,y]/\langle x^{2}-1,xy-y\rangle$ 的任意元素都可写为 $[ax+b+yg(y)]$, 其中 $a,b\in\mathbb{C}$ 且 $g(y)\in\mathbb{C}[y]$. 直接计算可知, 这样的元素 $[ax+b+yg(y)]$ 是多项式
\[
  f(z)=[x-1]z+[ax-bx+b-a]\in
  \bigl(\mathbb{C}[x]/\langle x^{2}-1\rangle\bigr)[z]
\]
的根, 这说明第二个扩张也是代数扩张.

然而, 虽然两个扩张都是代数扩张, 扩张
\[
  \mathbb{C}\subseteq\frac{\mathbb{C}[x,y]}{\langle x^{2}-1,xy-y\rangle}
\]
却不是代数扩张. 为说明这一点, 只要论证 $[y]$ 在 $\mathbb{C}$ 上不代数. 事实上, $\langle x^{2}-1,xy-y\rangle\subseteq\mathbb{C}[x,y]$ 的每个非零元素必是 $x-1$ 的倍式, 因此对任意非零多项式 $f(z)\in\mathbb{C}[z]$, 都不会有
\[
  f([y])=[f(y)]=0\in\frac{\mathbb{C}[x,y]}{\langle x^{2}-1,xy-y\rangle}.
\]
\end{Example}

虽然 Corollary~\ref{cor:algebraic-field-extensions-transitive} 对一般的环不成立, 但通过过渡到分式域, 它在整环这一特殊情形下仍然成立 (Exercise~\ref{ex:algebraic-transitivity-for-domains}).

本节最后举几个具体例子说明超越基的概念, 从空集这个例子一直到 $\mathbb{R}$ 中有关超越数的一个长期未解问题.

\begin{Example}{代数扩张的超越基}\leavevmode\label{ex:empty-transcendence-basis}
域扩张 $K\subseteq L$ 有空超越基, 当且仅当 $L$ 在 $K$ 上代数. 为验证这一点, 先注意 $\varnothing\subseteq L$ 自然是代数独立的: 由于 $\varnothing$ 不含任何元素, 元素之间就不可能存在代数关系. 因此, 按定义, $\varnothing\subseteq L$ 是超越基当且仅当 $L$ 在 $K(\varnothing)=K$ 上代数.
\end{Example}

\begin{Example}{单位球面的超越基}\leavevmode\label{ex:transcendence-basis-of-unit-sphere}
令 $K=\mathbb{C}$, 并令 $X=\mathcal{V}(x^{2}+y^{2}+z^{2}-1)\subseteq\mathbb{A}^{3}$. 来验证 $\{[x],[y]\}$ 是域扩张 $\mathbb{C}\subseteq\mathbb{C}(X)$ 的一个超越基. 为此要检查两个条件.

先验证 $\{[x],[y]\}$ 在 $\mathbb{C}$ 上代数独立. 设有多项式关系 $0=f([x],[y])$, 必须证明 $f=0$. 由于 $0=f([x],[y])=[f(x,y)]\in\mathbb{C}[X]$, 有
\[
  f(x,y)\in\mathcal{I}(X)=\langle x^{2}+y^{2}+z^{2}-1\rangle.
\]
然而, $\langle x^{2}+y^{2}+z^{2}-1\rangle$ 的每个非零元素关于 $z$ 都有正次数, 而 $f(x,y)$ 并非如此. 这意味着 $f(x,y)$ 必须是零多项式, 因此 $\{[x],[y]\}$ 在 $\mathbb{C}$ 上代数独立.

为证明 $\mathbb{C}(X)$ 在 $\mathbb{C}([x],[y])$ 上代数, 由 Proposition~\ref{prop:extending-fields-by-algebraic-elements} 及 $\mathbb{C}(X)=\mathbb{C}([x],[y],[z])$, 只需证明 $[z]$ 在 $\mathbb{C}([x],[y])$ 上代数. 这由 $[z]$ 是下列多项式的零点得到:
\[
  g(w)=w^{2}+([x]^{2}+[y]^{2}-1)\in\mathbb{C}([x],[y])[w].
\]
\end{Example}

\begin{Example}{$\mathcal{V}(f)$ 的超越基}\leavevmode\label{ex:transcendence-basis-of-hypersurface}
推广上一例, 若对某个不属于 $K[x_{1},\ldots,x_{n-1}]$ 的不可约多项式 $f\in K[x_{1},\ldots,x_{n}]$, 有 $X=\mathcal{V}(f)$, 则与上一例类似的论证表明 (Exercise~\ref{ex:hypersurface-transcendence-basis}), $\{[x_{1}],\ldots,[x_{n-1}]\}$ 是域扩张
\[
  K\subseteq K(X)=\operatorname{Frac}(K[x_{1},\ldots,x_{n}]/\langle f\rangle)
\]
的一个超越基.
\end{Example}

\begin{Example}{$\pi$ 和 $e$}\leavevmode\label{ex:pi-and-e-transcendence-basis}
判定 $\pi$ 和 $e$ 在 $\mathbb{Q}$ 上是否代数独立, 是一个长期未解的问题. 若它们代数独立, 则 $\{\pi,e\}$ 是 $\mathbb{Q}(\pi,e)$ 在 $\mathbb{Q}$ 上的一个超越基. 否则, 单元素集 $\{\pi\}$ 和 $\{e\}$ 各自都是 $\mathbb{Q}(\pi,e)$ 在 $\mathbb{Q}$ 上的一个超越基.
\end{Example}

看过几个超越基的例子后, 自然会有些疑问. 每个域扩张都有超越基吗? 如果有, 关于超越基的大小能说些什么? 下一节将证明, 若 $K\subseteq L$ 是有限生成域扩张, 则存在有限超越基, 且它们的大小都相同. 这样, 就能把这种域扩张的超越次数定义为任一超越基的大小, 类似于向量空间维数的一种定义方式. 超越次数这个概念是第~\ref{sec:dimension-of-affine-varieties}~节严格定义仿射簇维数的关键.

\BookExerciseHeading
\begin{BookExercises}
  \item\label{ex:generated-field-minimality} 设 $K\subseteq L$ 为域扩张, 且 $S\subseteq L$ 为子集.
  \begin{enumerate}
    \item 证明 $K(S)$ 是 $L$ 的子域.
    \item 若 $M$ 是 $L$ 中包含 $K$ 和 $S$ 的子域, 证明 $K(S)\subseteq M$.
  \end{enumerate}
  \item\label{ex:generated-field-in-stages} 设 $K\subseteq L$ 为域扩张, 且 $S_{1},S_{2}\subseteq L$ 为子集. 证明
  \[
    K(S_{1})(S_{2})=K(S_{1}\cup S_{2}).
  \]
  \item\label{ex:algebraic-over-generated-field-iff-algebraic-over-algebra} 设 $K\subseteq L$ 为域扩张, 且 $S\subseteq L$ 为集合. 证明 $L$ 在 $K(S)$ 上代数当且仅当 $L$ 在 $K[S]$ 上代数.
  \item\label{ex:algebraic-transitivity-for-domains} 设 $K$ 为域, 并设 $A\subseteq B$ 是有限生成 $K$-代数, 使得 $A$ 在 $K$ 上代数, 且 $B$ 在 $A$ 上代数. 假设 $A$ 和 $B$ 都是整环, 证明 $B$ 在 $K$ 上代数. (提示: 使用分式域.)
  \item\label{ex:hypersurface-transcendence-basis} 设 $f\in K[x_{1},\ldots,x_{n}]$ 是不属于 $K[x_{1},\ldots,x_{n-1}]$ 的不可约多项式, 考虑不可约仿射簇 $X=\mathcal{V}(f)$. 证明 $\{[x_{1}],\ldots,[x_{n-1}]\}$ 是 $K(X)$ 在 $K$ 上的一个超越基.
  \item\label{ex:noether-basis-is-transcendence-basis-of-fraction-field} 设 $A$ 是有限生成 $K$-代数, 且是整环. 证明 $A$ 在 $K$ 上的每个 Noether 基都是 $\operatorname{Frac}(A)$ 在 $K$ 上的一个超越基.
\end{BookExercises}
